%%=============================================================================
%% Keuzes technologieën
%%=============================================================================
\chapter{Keuzes technologieën}

\section{Blazor WebAssembly}

Voor de frontend werd gekozen voor Blazor WebAssembly. Dit is een framework 
van Microsoft waarmee je webapplicaties kan bouwen met C\# in plaats van 
JavaScript. Omdat Victor Buyck Steel Construction al ervaring heeft met het 
.NET-ecosysteem was dit een logische keuze.
\bigskip
Blazor WebAssembly draait volledig in de browser via WebAssembly, zonder 
dat er een constante serververbinding nodig is. De applicatie wordt één 
keer gedownload en kan daarna lokaal uitgevoerd worden. Dit is ideaal voor 
een onthaaltoestel dat de hele dag moet blijven werken.

\section{C\# en .NET}

De backend is geschreven in C\# en draait als een REST API. C\# werd gekozen 
omdat het aansluit bij de bestaande kennis binnen het bedrijf en goed 
integreert met Blazor. De frontend communiceert met de backend via HTTP-calls 
naar de API.

\section{MySQL en Entity Framework}

Voor de database werd gekozen voor MySQL. De data wordt beheerd via Entity 
Framework met de code-first aanpak. Dit betekent dat de databasestructuur 
wordt gegenereerd vanuit C\#-klassen, zonder dat er manueel SQL geschreven 
moet worden. Wijzigingen aan de database worden bijgehouden via migrations.

\section{Projectstructuur}

Het project is opgesplitst in vier aparte onderdelen. De \textbf{frontend} 
is het Blazor WebAssembly project dat de gebruikersinterface bevat. De 
\textbf{API} is de backend die de logica verwerkt en communiceert met de 
database. Het \textbf{DB}-project bevat de Entity Framework configuratie 
en de migrations. Ten slotte bevat het \textbf{DBEntities}-project de 
C\#-klassen die de databasetabellen voorstellen en die gedeeld worden 
tussen de API en het DB-project.
